\section{Taking shortcuts (lossy)：减少 KV cache 和模型字节}

前面的 workload 分析指出 generation 常被权重与 KV cache 读取限制，因此本部分允许改变模型表示来减少 bytes。GQA/MQA、MLA/CLA、local attention、quantization 与 pruning 都可能换来更高吞吐，但它们会改变可表达函数或数值精度，必须用任务质量和长上下文评测重新确认边界。

\subsection{Grouped-query attention (GQA)：按 head 共享 KV}

本节先从最直接的 KV head 共享开始。MHA 为每个 query head 保存独立 K/V，MQA 让所有 query heads 共享一组 K/V，GQA 则在两者之间按 group 共享；$N/K$ 越大，cache 越小、读取越少，但不同 heads 可使用的独立记忆子空间也越少。

\begin{figure}[H]
\centering
\includegraphics[width=0.9\textwidth]{gmqa.png}
\caption{Grouped-query attention：\(N\) 个 query heads 共享 \(K\) 个 key/value heads，\(K=1\) 是 MQA，\(K=N\) 是 MHA。}
\figsource{Scaling Book GQA/MQA diagram，本地化后供 CS336 Lecture 10 使用。}
\end{figure}

\begin{knowledgebox}{读图：GQA 图应该看 \(N/K\)}
MHA 为每个 query head 保存独立 KV；MQA 所有 query heads 共享一个 KV head；GQA 介于两者之间，每组 query heads 共享一组 KV。KV cache 大小约按 \(N/K\) 缩小，因此 latency 和 throughput 改善来自更少 HBM 读取。
\end{knowledgebox}

\begin{figure}[H]
\centering
\includegraphics[width=0.78\textwidth]{gqa-speed.png}
\caption{GQA 的 latency/throughput 结果：减少 KV heads 后更容易 fit，throughput 改善。}
\figsource{CS336 Lecture 10 官方讲义引用图。}
\end{figure}

\begin{knowledgebox}{读图：GQA speed 图说明什么}
横向比较时，KV cache 小的配置通常 latency 更低或能容纳更大 batch，从而 throughput 更好。若原 MHA 配置不 fit，GQA 首先解决容量；若能 fit，GQA 仍可能通过减少内存流量改善吞吐。
\end{knowledgebox}

\begin{figure}[H]
\centering
\includegraphics[width=0.85\textwidth]{gqa-accuracy.png}
\caption{GQA accuracy 检查：减少 KV cache 后必须验证准确率是否下降。}
\figsource{CS336 Lecture 10 官方讲义引用图。}
\end{figure}

\begin{warningbox}{GQA 是 lossy shortcut}
GQA 改变 attention 结构，通常会影响模型质量。速度图只说明系统收益，accuracy 图才说明是否值得用。课程主线是：有损 shortcut 必须以质量检查收尾。
\end{warningbox}

\begin{knowledgebox}{课堂提示：GQA 的速度来自 \(N/K\) 倍 KV 缩减}
老师先重申 inference 的瓶颈是 memory，再给出 GQA 的直接机制：\(N\) 个 query heads 只配 \(K\) 个 KV heads，cache 约缩小 \(N/K\) 倍，因此 latency/throughput 改善来自更少 HBM bytes。源码随后立即要求检查 accuracy，说明“更快且 fit”仍不等于可以无条件替换 MHA。
\end{knowledgebox}

\begin{knowledgebox}{讲义提醒：速度图和准确率图必须成对出现}
GQA 的价值不能只看 tokens/s。更少 KV heads 可能让更大 batch fit 进显存，所以性能收益包含容量效应；同时质量下降可能只在长上下文、知识密集或特定下游任务暴露。部署决策必须在目标 workload 上同时画出 latency/throughput 与 accuracy frontier。
\end{knowledgebox}

\subsection{Multi-head latent attention (MLA)：压缩 KV}

GQA 沿 head 维共享 K/V，MLA 则进一步改变 cache 的表示：保存低维 latent，并在使用时投影成 K/V。它减少长期存储和读取，却引入重建计算与潜在表达瓶颈，需要在相同 cache budget 下比较质量。

\begin{figure}[H]
\centering
\includegraphics[width=0.9\textwidth]{mla-schema.png}
\caption{Multi-head latent attention：存压缩 latent \(c=W_c h\)，需要时再投影出 \(K,V\)。}
\figsource{CS336 Lecture 10 官方讲义引用图。}
\end{figure}

\begin{knowledgebox}{读图：MLA 如何压缩 KV cache}
普通 attention 存 \(K=W_Kh,V=W_Vh\)，维度约为 \(N H\)。MLA 存低维 latent \(c\)，例如 DeepSeek v2 中从 \(16384\) 维压到 \(512\) 维，再额外保留 RoPE 所需维度。它把长期 cache 从高维 KV 变成低维 latent，换来每步投影计算。
\end{knowledgebox}

\begin{figure}[H]
\centering
\includegraphics[width=0.9\textwidth]{mla-accuracy.png}
\caption{MLA/MHA/GQA 准确率比较之一：先确认 MHA 与 GQA 的质量差异。}
\figsource{CS336 Lecture 10 官方讲义引用图。}
\end{figure}

\begin{figure}[H]
\centering
\includegraphics[width=0.8\textwidth]{mla-accuracy2.png}
\caption{MLA 准确率比较之二：MLA 在某些设置下可接近甚至超过 MHA，同时更省 KV cache。}
\figsource{CS336 Lecture 10 官方讲义引用图。}
\end{figure}

两张 accuracy 表需要作为同一组证据读取：先比较 MHA 与 GQA 的基线差异，再观察 MLA 在相同或更小 cache budget 下是否保持质量。若不同实验的模型大小、训练 token 或 cache 尺寸不一致，就不能把表中绝对分数直接解释成 MLA 机制本身的收益。

\begin{knowledgebox}{读图：MLA accuracy 两张表一起看}
第一张表确认不同 attention 结构的质量基线；第二张表展示 MLA 的质量和成本折中。读表时不能只看平均分，还要看任务类别和模型规模，因为 KV 压缩可能对长上下文、检索或数学任务影响不同。
\end{knowledgebox}

\subsection{Cross-layer attention (CLA)：跨层共享 KV}

MLA 压缩每一层保存的表示，CLA 则改变“哪些层必须各自拥有 cache”：相邻或分组 layers 共享 K/V，使层维度上的冗余像 GQA 在 head 维度上的冗余一样被压缩。代价是不同层无法自由保存各自的历史表示，因此应看 accuracy--cache Pareto frontier。

\begin{figure}[H]
\centering
\includegraphics[width=0.72\textwidth]{cla-diagram.png}
\caption{Cross-layer attention：跨层共享 KV，类似 GQA 在 heads 维共享 KV。}
\figsource{CS336 Lecture 10 官方讲义引用图。}
\end{figure}

\begin{figure}[H]
\centering
\includegraphics[width=0.85\textwidth]{cla-results.png}
\caption{CLA 结果：改善 accuracy 与 KV cache size 的 Pareto frontier。}
\figsource{CS336 Lecture 10 官方讲义引用图。}
\end{figure}

\begin{knowledgebox}{读图：CLA 的 Pareto frontier}
CLA 不是只追求最小 KV cache，而是在相同 cache budget 下尽量提高 accuracy，或在相同 accuracy 下减少 cache。Pareto frontier 向左下移动才说明结构改动有系统价值。
\end{knowledgebox}

\subsection{Local (sliding window) attention 与 hybrid/DeepSeek attention}

前面的压缩方法仍让每个新 token 访问全部历史 cache，本节转向选择“访问哪些位置”。局部注意力与 sliding-window attention 限制每层感受野，hybrid 模型周期性插入 full 或 global 层恢复远程信息，更新的压缩与选择机制则试图在超长上下文中只读取少量高价值状态。

\begin{figure}[H]
\centering
\includegraphics[width=0.9\textwidth]{longformer-attention.png}
\caption{Local/sliding-window attention：每层只看局部窗口，多层叠加后 effective context 线性扩大。}
\figsource{Longformer/local attention diagram，供 CS336 Lecture 10 使用。}
\end{figure}

\begin{knowledgebox}{读图：local attention 的收益和代价}
局部窗口让每层 KV cache 与可见窗口相关，而不是与完整 sequence length 线性相关。多层堆叠可以传播信息，但远距离精确依赖会变弱。因此实践中常用 hybrid layers：部分层 local，部分层 global。
\end{knowledgebox}

\begin{figure}[H]
\centering
\includegraphics[width=0.9\textwidth]{deepseek-v4-attention.png}
\caption{DeepSeek v4 attention：CSA、DSA、HCA 等多级压缩/选择支持超长上下文。}
\figsource{CS336 Lecture 10 官方讲义引用图。}
\end{figure}

这张结构图应沿数据流阅读：先看历史状态如何被压缩成候选表示，再看 selector 如何减少需要精确 attention 的位置，最后检查被选状态如何进入当前 token 计算。它表达的是“先便宜筛选、再昂贵计算”的分层机制，而不是免费获得无限上下文；selector 误差和额外索引开销仍需计入。

\begin{knowledgebox}{术语消化：DeepSeek v4 attention 三个缩写}
\begin{tabular}{L{0.18\textwidth}L{0.34\textwidth}L{0.38\textwidth}}
\toprule
术语 & 机制 & 推理意义 \\
\midrule
CSA & Compressed Sparse Attention，把每 \(m\) 个 tokens 压成一个表示。 & 降低长上下文 KV 读取。 \\
DSA & DeepSeek Sparse Attention，选择 top \(k\) 相关位置。 & 把注意力集中到更有用 token。 \\
HCA & Heavily Compressed Attention，更激进压缩。 & 服务超长上下文时进一步控成本。 \\
\bottomrule
\end{tabular}
\end{knowledgebox}

\subsection{quantization 与 Activation-aware quantization (AWQ)}

Attention 结构优化之后，本节处理另一大类 bytes：模型权重与 activation。Quantization 用更少 bit 表示数值，直接降低 HBM 容量和带宽压力，并可能启用更快低精度矩阵单元；代价是 rounding、clipping 和 outlier 误差，因此训练时量化与训练后量化有不同适用边界。

\begin{figure}[H]
\centering
\includegraphics[width=0.75\textwidth]{fp8-quantization.png}
\caption{不同数值精度示意：fp32、bf16、fp8、int8、int4 代表不同 bytes 和误差。}
\figsource{CS336 Lecture 10 官方讲义引用/本地化图。}
\end{figure}

读这张格式图时要同时看 exponent、mantissa 与总 bytes：FP8 相比 BF16/FP16 降低存储和带宽，但动态范围与精度由具体 E4M3/E5M2 变体决定；INT8/INT4 还需要 scale、zero point 或 group-wise metadata。理论上的 2 倍/4 倍压缩不会自动转化成同等速度，kernel 和 dequantization 也必须匹配。

\begin{knowledgebox}{读图：量化为什么直接影响推理}
推理 memory-bound 时，参数和 KV cache 的 bytes 直接进入 latency/throughput。bf16 是常见默认；fp8/int8/int4 减少内存流量，但会引入量化误差。位宽越低，越需要校准、补偿或训练时适配。
\end{knowledgebox}

\begin{knowledgebox}{术语消化：QAT、PTQ、GPTQ、AWQ}
\begin{tabular}{L{0.18\textwidth}L{0.34\textwidth}L{0.38\textwidth}}
\toprule
方法 & 做法 & 代价/适用性 \\
\midrule
QAT & Quantization-aware training，训练中模拟量化误差。 & 质量好但昂贵。 \\
PTQ & Post-training quantization，训练后用校准数据确定 scale/zero point。 & 便宜，常用于部署。 \\
GPTQ & 用 Hessian 信息补偿量化误差。 & 更精细但实现复杂。 \\
AWQ & Activation-aware quantization，保留少量重要权重高精度。 & 关注大 activation channel 命中的权重。 \\
\bottomrule
\end{tabular}
\end{knowledgebox}

\begin{figure}[H]
\centering
\includegraphics[width=0.95\textwidth]{awq-schema.png}
\caption{AWQ：根据 activation 重要性选择少量权重保留更高精度。}
\figsource{CS336 Lecture 10 官方讲义引用图。}
\end{figure}

AWQ 的关键不是把所有权重一视同仁地压到低 bit，而是根据 activation 统计识别对输出更敏感的通道，并通过 scaling 或少量高精度保留降低误差。读图时应区分离线校准、权重变换、量化存储和运行时反量化四个阶段，避免把校准成本误算成每 token 开销。

\begin{knowledgebox}{读图：AWQ schema 说明什么}
AWQ 的直觉是并非所有权重量化误差同等重要。若某些 activation channels 很大，与其相乘的权重误差会被放大。保留 0.1\%-1\% 重要权重的高精度，可能用很小内存代价换来明显质量改善。
\end{knowledgebox}

\begin{knowledgebox}{课堂提示：AWQ 让 activation 决定哪些权重值得精度}
老师把 AWQ 的观察压缩成三步：某些 activation channels 很大；与它们相乘的权重更重要；因此只为约 0.1\%--1\% 的关键权重分配更多精度。课程给出的案例从 fp16 到 int3 可带来约 4 倍更低内存和 3.2 倍 speedup，但这些数字依赖校准数据、kernel 与目标模型，机制比单一 benchmark 更可迁移。
\end{knowledgebox}

\subsection{pruning 与 distillation}

量化保留稠密结构但降低每个元素的 bit，本节则直接删除权重、通道、层或注意力头。Pruning 会改变模型容量和 kernel shape，若硬件无法利用非结构化稀疏，参数更少也未必更快；distillation 用原模型输出修复裁剪后的行为，因此两者应作为一个闭环而不是两个独立技巧。

\begin{figure}[H]
\centering
\includegraphics[width=0.9\textwidth]{pruning-kd-loop.png}
\caption{Pruning + distillation loop：识别重要结构、裁剪，再用原模型蒸馏修复。}
\figsource{CS336 Lecture 10 官方讲义引用图。}
\end{figure}

\begin{knowledgebox}{读图：pruning 为什么必须 distill}
直接移除 layers、heads 或 hidden dimensions 会破坏模型函数。蒸馏让小模型模仿原模型输出，把结构损伤修复回来。它是有损 shortcut，但通过 teacher signal 降低质量损失。
\end{knowledgebox}

\begin{knowledgebox}{老师强调：先 rip out，再 fix it up}
源码用非常直接的语言概括 pruning：先从昂贵模型中移除不重要的 layer、head 或 hidden dimension，再用原模型蒸馏修复。它同时对应两种 recipe：from scratch 是定义并训练更快架构；distillation 则从原模型初始化不同结构，再通过 teacher signal repair。少参数只是开始，修复后的质量才决定 shortcut 是否成立。
\end{knowledgebox}

\begin{figure}[H]
\centering
\includegraphics[width=0.78\textwidth]{pruning-kd.png}
\caption{Pruning + KD 结果：裁剪后经蒸馏可恢复部分质量，同时降低推理成本。}
\figsource{CS336 Lecture 10 官方讲义引用图。}
\end{figure}

\begin{warningbox}{有损 shortcut 的共同风险}
GQA、MLA、CLA、local attention、quantization、pruning 都可能改变输出分布。它们的验收不能只看速度，必须看准确率、长上下文能力、校准、下游任务和安全性指标。
\end{warningbox}

